% ============================================================
\chapter{Nhật ký thực hiện}
\label{app:nhatky}

Bảng~\ref{tab:nhatky} ghi lại nhật ký thực hiện đồ án theo từng tuần, bao gồm nội dung công việc đã hoàn thành, kết quả đạt được và nhận xét của Giảng viên hướng dẫn trong các buổi gặp định kỳ.

\renewcommand{\arraystretch}{1.3}
\begin{longtable}{|c|p{2.0cm}|p{2.0cm}|p{5.0cm}|p{3.0cm}|}
    \caption{Nhật ký thực hiện đồ án}
    \label{tab:nhatky} \\
    \hline
    \textbf{Tuần} & \textbf{Từ ngày} & \textbf{Đến ngày} & \textbf{Tóm tắt nội dung} & \textbf{Nhận xét GVHD} \\ \hline
    \endfirsthead
    \multicolumn{5}{r}{\footnotesize\textit{(tiếp theo)}} \\
    \hline
    \textbf{Tuần} & \textbf{Từ ngày} & \textbf{Đến ngày} & \textbf{Tóm tắt nội dung} & \textbf{Nhận xét GVHD} \\ \hline
    \endhead
    \hline \multicolumn{5}{r}{\footnotesize\textit{(còn tiếp)}} \\
    \endfoot
    \hline
    \endlastfoot
    1 & 20/01/2026 & 25/01/2026 & Tìm hiểu bài toán, nghiên cứu các nền tảng IELTS hiện có, xác định hướng đề tài & Đồng ý hướng đề tài, yêu cầu làm rõ phạm vi chức năng \\ \hline
    2 & 26/01/2026 & 01/02/2026 & Phân tích nhu cầu người dùng, viết tài liệu đặc tả yêu cầu sơ bộ & Góp ý bổ sung tính năng Mock Test đầy đủ 4 kỹ năng \\ \hline
    3 & 02/02/2026 & 08/02/2026 & Nghiên cứu NestJS, Prisma, PostgreSQL, RabbitMQ, Gemini API, faster-whisper; lập bảng so sánh công nghệ & Nhất trí tech stack, yêu cầu giải thích rõ lý do chọn NestJS \\ \hline
    4 & 09/02/2026 & 15/02/2026 & Vẽ use-case diagram, thiết kế ERD PostgreSQL (14 bảng), sơ đồ lớp & Cần bổ sung thêm Alternative Flow trong đặc tả UC \\ \hline
    5 & 16/02/2026 & 22/02/2026 & Thiết kế wireframe Web và Mobile, xây dựng sơ đồ luồng màn hình & Giao diện khá trực quan, cần thêm màn hình tiến độ học tập \\ \hline
    6 & 23/02/2026 & 01/03/2026 & Khởi tạo dự án NestJS Modular Monolith, cấu hình PostgreSQL, RabbitMQ và Prisma schema & Kiến trúc module hoá hợp lý, cần chú ý API versioning \\ \hline
    7 & 02/03/2026 & 08/03/2026 & Hoàn thiện User Service, Content Service; viết API đăng ký, đăng nhập, JWT & API hoạt động tốt, bổ sung cơ chế refresh token \\ \hline
    8 & 09/03/2026 & 15/03/2026 & Lập trình giao diện Next.js: Landing page, dashboard, module Vocabulary & Giao diện đẹp, cần tối ưu responsive trên màn hình nhỏ \\ \hline
    9 & 16/03/2026 & 22/03/2026 & Hoàn thiện Web: Listening, Reading, Writing (tích hợp Gemini API scoring) & AI scoring Writing hoạt động ổn, cần xử lý trường hợp AI timeout \\ \hline
    10 & 23/03/2026 & 29/03/2026 & Khởi tạo dự án React Native, xây dựng màn hình đăng nhập và từ vựng & Ứng dụng chạy tốt trên Android, tiếp tục hoàn thiện \\ \hline
    11 & 30/03/2026 & 05/04/2026 & Hoàn thiện Mobile: Listening, Speaking (faster-whisper), Mock Test & STT hoạt động tốt, cần cải thiện UI feedback khi đang ghi âm \\ \hline
    12 & 06/04/2026 & 12/04/2026 & Tích hợp AI Conversation Speaking, tinh chỉnh prompt engineering & AI hội thoại tự nhiên, cần giới hạn số lượt gọi API/người dùng \\ \hline
    13 & 13/04/2026 & 19/04/2026 & Thực hiện 22 test case, ghi lại kết quả; sửa 4 lỗi phát hiện trong kiểm thử & Kết quả kiểm thử tốt, cần bổ sung test case cho edge case \\ \hline
    14 & 20/04/2026 & 26/04/2026 & Deploy lên GCP (Cloud Run cho frontend, GCP VM + Docker Compose cho backend), thu thập phản hồi từ 5 người dùng thử & Sản phẩm hoàn thiện, tập trung viết báo cáo \\ \hline
    15 & 27/04/2026 & 05/05/2026 & Hoàn thiện báo cáo đồ án, chuẩn bị slide thuyết trình, nộp báo cáo & Báo cáo đạt yêu cầu, duyệt nộp \\ \hline
\end{longtable}

